\chapter{重整化攻坚战}

式~(\ref{eq:quasipdf}) 中的 Wilson 线是自能发散的对象：除通常的对数
外，它携带正比于其长度 $z$ 的\emph{线性} UV 发散。能否以及如何驯服它，
决定了格点部分子物理的命运。本章以七篇论文追踪这场战役：从让格点
矩阵元赢得尊重的老牌重整化方案，到可重整性的证明，再到本库日常使用
的实用方案（比值/RI/MOM 混合、辅助场、自重整化）。

% ----------------------------------------------------------------------
\section{RI/MOM：老牌方案}
\papercard{G.\ Martinelli, C.\ Pittori, C.\ T.\ Sachrajda, M.\ Testa,
A.\ Vladikas}
{格点算符非微扰重整化的一般方法}
{Nucl.\ Phys.\ B \textbf{445}, 81--108 (1995)}
{arXiv:hep-lat/9411010 | DOI: 10.1016/0550-3213(95)00126-D}
\whyselected{非微扰算符重整化的奠基方案（显式引用约 $\sim$8 次，加上
普遍的隐式使用）：库内每个 RI/MOM 或 SMOM 重整化的矩阵元皆源自于此。}
\physics{在一族 Landau 规范离壳动量减除态上对矩阵元施加重整化条件来
重整算符 $O$：
\begin{equation}
  Z_O(\mu,a)\,\Lambda_O(p)\big|_{p^2=\mu^2} = \Lambda_O^{tree}(p),
  \qquad p^2 \gg \Lambda_{\rm QCD}^2 ,
\end{equation}
由无微扰的模拟计算。该方案重整化群封闭、具备到 $\overline{\rm MS}$
的可计算转换因子；而对后续章节至关重要的是——它可以逐字推广到双局域
算符，其动量空间振幅顶点成为诊断 Wilson 线线性发散的仪器。}
\impact{库内所有早期准 PDF 提取（Liu 等、物理点螺旋度、同位旋矢量
系统研究）都以 RI/MOM 风格重整化；TMD 论文记录了它失效之处（第六章
语境），而这本身也是本文遗产的一部分：正因为方案存在，才能诊断其
边界。}
\cardend

% ----------------------------------------------------------------------
\section{可重整性之证明}
\papercard{T.\ Ishikawa, Y.-Q.\ Ma, J.-W.\ Qiu, S.\ Yoshida}
{准部分子分布函数的可重整性}
{Phys.\ Rev.\ D \textbf{96}, 094019 (2017)}
{arXiv:1707.03107 | DOI: 10.1103/PhysRevD.96.094019}
\whyselected{准 PDF 算符可重整性的第一个一般证明（约 $\sim$13/99）：
没有它，整个纲领可能沦为数值假象。}
\physics{分析双局域算符的全部单圈图后，论文证明所有幂发散都定域于
Wilson 线自能的非 cusp 部分——即算符在等维度同量子数的算符之间可乘
地重整化（至平凡混合），无需新的非局域反项。结合规范不变的 UV
调节器论证，结果保证恰当重整化的准 PDF 存在连续极限。Ma--Qiu 的姊妹
PRL（arXiv:1404.6860 谱系；"good lattice cross sections"）把同一逻辑
框定为一般策略：在格点可靠的短距离处计算，再匹配到长距离可观测量。}
\impact{库内的每一个重整化提案——比值方案、辅助场减除、混合方案——
都引用这一证明作为理论担保。}
\cardend

% ----------------------------------------------------------------------
\section{驯服线性发散}
\papercard{J.-W.\ Chen, X.\ Ji, J.-H.\ Zhang}
{通过 Wilson 线重整化改进的准部分子分布}
{Nucl.\ Phys.\ B \textbf{915}, 1--9 (2017)}
{arXiv:1609.08102 | DOI: 10.1016/j.nuclphysb.2016.12.011}
\whyselected{辅助场（质量反项）消除线性发散方案的源头（约
$\sim$11/99）：语料库中几乎每篇重整化论文都引用它。}
\physics{引入二次项携带质量 $m_0$ 的辅助场 $\phi$；开 Wilson 线在一个
扩大理论中的期望值里重新表述，其线性发散仅作为 $\phi$ 的自能出现，
被下式吸收：
\begin{equation}
  m_0^2(z,a)=m^2 + \frac{c_1}{a} + \cdots .
\end{equation}
减除后的改进准分布在微扰论的任意阶只含对数发散，从而接受标准的
非微扰重整化。论文还计算了首个格点正规化的单圈匹配核，暴露离散化
伪迹如何进入 $K$。}
\impact{这一构造是 Green--Jansen--Steffens 实现与下文混合方案的种子；
概念上它把"减除一个线性发散"从数值问题转化为质量重整化问题——可以
通过跨 $z$ 拟合 $m_0(z)$ 解决。}
\cardend

% ----------------------------------------------------------------------
\section{重整化纲领的形式化}
\papercard{X.\ Ji, J.-H.\ Zhang, Y.\ Zhao}
{大动量有效理论中部分子物理的重整化}
{Phys.\ Rev.\ Lett.\ \textbf{120}, 112001 (2018)}
{arXiv:1706.08962 | DOI: 10.1103/PhysRevLett.120.112001}
\whyselected{LaMET 重整化的纲领性陈述（约 $\sim$13/99）：辅助场形式
体系及进入 NNLO 匹配的 Wilson 线反常量纲的标准引用。}
\physics{该快报系统化了如今处处使用的两步逻辑：(i)~在选定方案内
（RI/MOM、比值或辅助场质量）非微扰地重整准算符，使矩阵元在连续极限
下有限；(ii)~转换到 $\overline{\rm MS}$ 并匹配到光锥分布，核中包含
cusp 对数 $\ln(P^z/\mu)$。文章还计算了单圈转换因子，并按阶演示最终
光锥 PDF 的方案无关性。}
\impact{库内的混合重整化应用（横率、Boer--Mulders、胶子 PDF）以此
快报为定义框架；其记号是 LP3 论文事实上的标准。}
\cardend

% ----------------------------------------------------------------------
\section{辅助场的实现}
\papercard{J.\ Green, K.\ Jansen, F.\ Steffens}
{由格点 QCD 得到非局域夸克双线性的非微扰重整化}
{Phys.\ Rev.\ Lett.\ \textbf{121}, 022004 (2018)}
{arXiv:1707.07152 | DOI: 10.1103/PhysRevLett.121.022004}
\whyselected{首次数值证明辅助场方案在真实系综上可行（约 $\sim$14/99）：
Chen--Ji--Zhang 思想与日常实践之间的桥梁。}
\physics{在 $N_f=2+1$ domain wall 系综上测量准 PDF 算符格林函数，跨
分离距离与格距拟合线性发散的 $\phi$ 场质量 $m_0(z,a)$，将其减除后展示
残余矩阵元按对数趋于有限连续极限——随后提供到 $\overline{\rm MS}$ 的
转换因子。这种对幂发散量的"拟合提取"及其关联系统误差处理，成为后来
自重整化（见下）推广的模板。}
\impact{在本库内，ETMC 与 LPC 路线为混合短距处理辩护时引用此文；其
误差分析模式重现于每一次现代提取。}
\cardend

% ----------------------------------------------------------------------
\section{完整的处方}
\papercard{C.\ Alexandrou, K.\ Cichy, M.\ Constantinou,
K.\ Hadjiyiannakou, K.\ Jansen, H.\ Panagopoulos, F.\ Steffens}
{准 PDF 的完整非微扰重整化处方}
{Nucl.\ Phys.\ B \textbf{923}, 394--415 (2017)}
{arXiv:1706.00265 | DOI: 10.1016/j.nuclphysb.2017.08.010}
\whyselected{单一方案最完备的处理（约 $\sim$10/99）：ETMC 学派的
标准，库内多篇论文（含 twisted mass 胶子 PDF）共享。}
\physics{对非极化、螺旋度与横率准 PDF，论文构造了完整的 RI$'$ 处方：
所有相关张量结构的重整化条件、经 $z$ 依赖质量项减除线性发散、纳入
twist-3 结构中标量算符的混合、以及单圈微扰转换因子
$\mathrm{RI}'\to\overline{\rm MS}$。这是一份模拟可以从头执行到尾的
菜谱——矩阵元、$Z$ 因子、匹配、外推——这解释了它作为领域实用基准的
地位。}
\impact{在本库内它支撑 ETMC 质子胶子 PDF 计算，并充当与 hybrid-ratio
方案对比的参照；规范固定精度论文恰好量化了其 Landau 规范步骤需要
多小心。}
\cardend

% ----------------------------------------------------------------------
\section{胶子算符的可乘重整化}
\papercard{Z.-Y.\ Li, Y.-Q.\ Ma, J.-W.\ Qiu}
{定义准部分子分布的算符之可乘重整化}
{Phys.\ Rev.\ Lett.\ \textbf{122}, 062002 (2019)}
{arXiv:1809.01836 | DOI: 10.1103/PhysRevLett.122.062002}
\whyselected{胶子准 PDF 的严格基础（约 $\sim$8/99）：库内几乎每篇
胶子 PDF 文章都引用。}
\physics{对带 Wilson 线的全部规范不变伴随双局域算符分类，并证明——
使用规范不变 UV 调节器使出现的只有真正的 UV 奇异性——全部 36 个
胶子准部分子算符的紫外发散定域化且\emph{可乘地}重整化：每个算符映到
对称性固定的有限组合，超出平凡情形无算符混合。这把 Ishikawa 等对
夸克的证明补全到了胶子，认证了格点胶子准 PDF 定义连续量。}
\impact{与 Zhang 等人的匹配分析（下一章）一道，这条定理就是"首个格点
胶子 PDF"论文得以存在的原因。库内胶子准 PDF 路线在每个方法论节点
引用它。}
\cardend
